\documentclass[11pt]{report}
\usepackage[a4paper,margin=1in]{geometry}
\usepackage{amsmath,amssymb}
\usepackage{booktabs}
\usepackage{longtable}
\usepackage{array}
\usepackage[colorlinks=true,linkcolor=blue,urlcolor=blue]{hyperref}
\usepackage{fancyhdr}
\usepackage{xcolor}
\usepackage{titlesec}
\usepackage{enumitem}

\definecolor{hbblue}{HTML}{2A4D8F}
\definecolor{warnred}{HTML}{9C2B2B}
\definecolor{boxgrey}{HTML}{F4F4F6}

\titleformat{\chapter}[display]{\bfseries\Large\color{hbblue}}{\chaptertitlename\ \thechapter}{10pt}{\huge}
\pagestyle{fancy}
\fancyhf{}
\fancyhead[R]{\small\itshape What Feeds the Acuity Number, and How Much}
\fancyfoot[C]{\thepage}

\newsavebox{\kbsave}
\newsavebox{\wbsave}
\newenvironment{keybox}[1]%
{\par\vspace{6pt}\begin{lrbox}{\kbsave}\begin{minipage}{0.88\textwidth}%
{\bfseries\color{hbblue}#1}\par\vspace{4pt}\ignorespaces}%
{\end{minipage}\end{lrbox}%
\begin{center}\setlength{\fboxrule}{1.2pt}\setlength{\fboxsep}{7pt}%
\fcolorbox{hbblue}{boxgrey}{\usebox{\kbsave}}\end{center}\vspace{6pt}\par}
\newenvironment{warnbox}[1]%
{\par\vspace{6pt}\begin{lrbox}{\wbsave}\begin{minipage}{0.88\textwidth}%
{\bfseries\color{warnred}#1}\par\vspace{4pt}\ignorespaces}%
{\end{minipage}\end{lrbox}%
\begin{center}\setlength{\fboxrule}{1.2pt}\setlength{\fboxsep}{7pt}%
\fcolorbox{warnred}{white}{\usebox{\wbsave}}\end{center}\vspace{6pt}\par}

\title{\textbf{\Huge What Feeds the Acuity Number}\\[0.3cm]
\textcolor{hbblue}{\huge Inputs, Weights, and Why Each One Is There}\\[0.5cm]
\large\itshape Every characteristic used by \texttt{predict\_acuity()}, its numerical
weight in the final result, and the source or reasoning behind it}
\author{}
\date{Traced against a live run: 400\,\textmu m, $M=8$\,D, verified 2026-09-09}

\begin{document}
\maketitle
\tableofcontents

\chapter{The final equation, and what each piece contributes}

\section{The formula}

\begin{equation}
\text{MAR}_{\text{arcmin}} = \underbrace{\frac{30}{f_{\text{MTF10}}}}_{\text{(1) optical cutoff}}
\times \underbrace{\text{VSOTF}^{-0.1}}_{\text{(2) neural weighting}},
\qquad
\text{logMAR} = \log_{10}(\text{MAR}) + \underbrace{\Delta_{\text{bright}}}_{\text{(3) illuminance}}
\end{equation}

Three characteristics of the optical system feed this equation. Nothing else does. This chapter
quantifies exactly how much each one moves the final number, using a live traced run at
400\,\textmu m, $M=8$\,D as the worked example throughout --- every number below was read directly
from a single call to \texttt{run\_stage4\_point()}, not estimated.

\begin{table}[h!]
\centering
\begin{tabular}{p{3.7cm}p{4.2cm}p{2.0cm}p{4.0cm}}
\toprule
\textbf{Characteristic} & \textbf{What it measures} & \textbf{Weight in final logMAR} &
\textbf{At this operating point} \\
\midrule
(1) MTF10 cutoff frequency & Optical resolving power of the pinhole+eye system & \textbf{dominant} &
14.44 cpd $\to$ logMAR $+0.3175$ \\
(2) VSOTF (neural-weighted Strehl) & How much of that optical signal survives visual processing &
\textbf{small} & 0.739 $\to$ shifts logMAR by $+0.0131$ \\
(3) Retinal illuminance & Whether the image is bright enough for full acuity & \textbf{usually zero}
& 1257\,Td, above the 300\,Td knee $\to$ shifts logMAR by $0.0000$ \\
\bottomrule
\end{tabular}
\caption{The three inputs, ranked by actual effect size at the design's own operating point.}
\end{table}

\begin{keybox}{The one-sentence summary}
Acuity is set almost entirely by where the modulation transfer function crosses 10\% contrast. The
neural weighting nudges it by about a hundredth of a logMAR unit. The brightness term is fully wired
in, cited, and currently contributes nothing to any of the eight fabricated designs --- because
every one of them, at the assumed scene luminance, sits above the threshold where it would matter.
\end{keybox}

\chapter{External validation status --- read this before quoting any predicted logMAR}

\begin{warnbox}{The predictor has not been validated against real clinical data, and the one attempt failed}
The full model (MTF10 + VSOTF + brightness correction, exactly as traced in this document) was run
against Wang et al.\ (2009)'s real clinical pinhole-acuity measurements, in
\texttt{stage4\_human\_vision/RESULTS.md}.

\begin{center}
\begin{tabular}{cccc}
\toprule
$M$ (D) & $R^2$ & Pearson $r$ & $p$-value \\
\midrule
0 & $-6.39$ & 0.395 & 0.510 \\
2 & $-4.98$ & 0.393 & 0.513 \\
5 & $-2.84$ & 0.405 & 0.499 \\
\bottomrule
\end{tabular}
\end{center}

Negative $R^2$ means the model predicts worse than simply guessing the mean every time. $p\approx0.50$
means no statistically significant relationship was detected. This is not weak correlation --- every
one of the 15 tested (aperture, $M$) points shows the \emph{same-direction} bias: the model
over-predicts blur (worse logMAR than observed) at every point tested.
\end{warnbox}

\textbf{The luminance assumption was tested and ruled out as the cause.} Correcting the scene
luminance from the 10\,000\,cd/m$^2$ default to Wang 2009's own reported illumination
(15 footcandles $\to$ 43.7\,cd/m$^2$) made $R^2$ \emph{worse} ($-7.88$ vs.\ $-6.68$ at $M=0$), because
it pushed apertures further below the 300\,Td knee and added predicted blur on top of a model that was
already over-predicting it (\texttt{LUMINANCE\_RECONCILIATION.md}).

\textbf{The project's own leading diagnosis, most likely cause first:} (1) the acuity formula has no
fitted offset or scale against real data --- \texttt{DESIGN.md} flags this as a known
``requires empirical calibration'' gap; (2) the Wang 2009 dataset is non-monotonic at $n=5$
(observed logMAR 0.28, 0.31, \textbf{1.00}, 0.40, 0.50 at 500/750/1000/1200/1500\,\textmu m --- a spike
at 1000\,\textmu m that reverses at 1200) and a single outlier at $n=5$ sinks any smooth model
regardless of which optical metric underlies it; (3) the dataset records no per-subject refractive
error, adding irreducible noise; (4) $n=5$ is too small for significance even with real signal
present. \texttt{RESULTS.md}'s own conclusion: this ``cannot be properly validated with currently
available public data.''

\begin{keybox}{How to use every number in this document, given the above}
Every logMAR value quoted here is a \textbf{self-consistent model output} --- traceable, cited term
by term, internally coherent --- \emph{not} a validated clinical prediction. Say the validation status
yourself if the acuity model comes up: internally consistent, externally unvalidated, the one
validation attempt failed, luminance was tested and ruled out as the cause, the leading suspects are
lack of calibration and a small non-monotonic clinical dataset rather than the underlying optics. This
is also the reason the project's own logMAR measurement protocol matters more than a confirmatory
afterthought --- it is calibration data this model currently has none of.
\end{keybox}

\chapter{Characteristic 1: MTF10 cutoff frequency --- the dominant term}

\section{What it is}

The spatial frequency, in cycles per degree, at which the eye's modulation transfer function falls to
10\% contrast. Below that frequency, detail is still transmitted with enough contrast to see; above
it, the eye cannot reliably tell light from dark at that fine a spacing.

\section{Where it comes from}

Computed in \texttt{optical\_quality.py} from the retinal point-spread function that Stage~2's wave
propagation produces. The PSF is radially averaged into an MTF via
\texttt{stage3\_metrics.mtf\_radial}, then the frequency at 10\% is read off by
\texttt{mtf50\_mtf10}. This is a direct measurement of the simulated optical system --- not a fitted
or assumed number.

\section{How it enters the acuity formula}

\begin{equation}
\text{MAR}_{10} = \frac{30}{f_{\text{MTF10}}\,[\text{cpd}]}
\end{equation}

The constant 30 is \texttt{cycles\_per\_degree\_at\_20\_20} --- the standard clinical definition that
20/20 (logMAR 0) acuity corresponds to resolving 30 cycles per degree. This is the entire
definitional anchor of the whole model: everything downstream is expressed relative to this one
reference number.

\textbf{Worked example, 400\,\textmu m at $M=8$\,D:} MTF10 $= 14.4431$\,cpd, so
$\text{MAR}_{10} = 30/14.4431 = 2.0771$ arcmin, giving $\log_{10}(2.0771) = +0.3175$ logMAR before any
other correction is applied. This is 96\% of the final answer ($+0.3306$).

\begin{keybox}{Why this is the dominant term by design, not by accident}
The model's philosophy is: acuity is fundamentally an optical-resolution question, and everything
else (neural weighting, illuminance) is a secondary correction on top of that resolution limit. That
is a defensible simplification --- MTF10 alone is a standard, clinically-used proxy for acuity in its
own right (it is how many autorefractors and wavefront-aberrometer-based acuity predictions work) ---
but it does mean the model is only as good as the PSF that produced this one number. Every validation
concern raised elsewhere in this project about the retinal PSF (the FWHM ring-hopping, the 8\% Airy
tolerance at the first null) is inherited directly into this term, because it is computed from the
same PSF.
\end{keybox}

\chapter{Characteristic 2: VSOTF --- neural-weighted optical quality}

\section{What it is}

Visual Strehl ratio computed from the OTF, or VSOTF: a single number summarising how much of the
optical signal survives once the eye's own contrast sensitivity is factored in, normalised against
the diffraction-limited case.

\section{Where it comes from}

\begin{equation}
\text{VSOTF} = \frac{\int \text{OTF}_{\text{real}}(f)\,\text{CSF}(f)\,df}
{\int \text{OTF}_{\text{real,diffraction-limited}}(f)\,\text{CSF}(f)\,df}
\end{equation}

computed in \texttt{visual\_strehl.py}. The numerator integrates the \emph{actual} optical transfer
function --- weighted by the eye's contrast sensitivity function (CSF), so frequencies the eye barely
perceives contribute little to the integral. The denominator is the same integral for a perfect,
diffraction-limited system at the same aperture. VSOTF $=1$ means the system performs as well as
physically possible; lower values mean real aberrations (here, defocus) are costing performance.

The CSF weighting uses the \textbf{Barten model} by default (\texttt{csf.py},
\texttt{barten\_simplified\_csf}), a luminance- and field-size-dependent closed form. Two alternative
CSF models (Mannos--Sakrison, a normalised Watson template) are computed alongside it for
comparison but are not used in the recommended acuity number.

\section{How it enters the acuity formula}

\begin{equation}
\text{MAR}_{\text{adjusted}} = \frac{\text{MAR}_{10}}{\text{clamp}(\text{VSOTF},\,0.25,\,2.0)^{0.1}}
\end{equation}

\textbf{Worked example:} VSOTF $=0.7389$. $0.7389^{0.1} = 0.9700$. $\text{MAR}_{\text{adjusted}} =
2.0771 / 0.9700 = 2.1409$ arcmin, moving logMAR from $+0.3175$ to $+0.3306$ --- a shift of
\textbf{0.0131 logMAR}.

\begin{keybox}{Why the exponent is 0.1, and why that makes this a deliberately weak correction}
The clamp restricts VSOTF to $[0.25, 2.0]$ before the exponent is applied, and the exponent itself is
small. Even at the clamp's lower bound, $0.25^{-0.1} = 1.149$, a MAR increase of only 15\%
($\approx 0.06$ logMAR). This is a design choice: VSOTF is being used as a \emph{fine correction} to
the MTF10-based estimate, not as an independent, comparably-weighted acuity predictor in its own
right. The reasoning is that MTF10 already captures most of what the degraded optics do to acuity;
VSOTF adds the additional information that some of the remaining contrast is more or less usable to
the visual system depending on which spatial frequencies survived, without being allowed to swing the
answer by more than a fraction of a line.
\end{keybox}

\section{What the Barten CSF weighting itself depends on}

Two further characteristics enter here, both fixed defaults rather than swept design variables in
your current results:

\begin{longtable}{p{3.5cm}p{2.5cm}p{7.5cm}}
\toprule
\textbf{Input} & \textbf{Default} & \textbf{Effect} \\
\midrule
\endhead
Scene luminance & 10\,000\,cd/m$^2$ & Higher luminance raises peak CSF sensitivity and shifts the
CSF's high-frequency cutoff outward (via the $a$ and $b$ terms in \texttt{barten\_simplified\_csf}) ---
brighter scenes let the eye use finer detail, so VSOTF's frequency weighting shifts toward higher
spatial frequencies at high luminance. \\
Field diameter & 20$^\circ$ & Enters the same $a$ term; a larger field lowers the $12/w$ term in the
denominator, slightly raising sensitivity. This models spatial summation across a wider retinal area.
\\
\bottomrule
\end{longtable}

\chapter{Characteristic 3: retinal illuminance --- present, cited, usually zero}

\section{What it is}

The amount of light actually reaching the retina, in trolands (Td) --- a unit that already accounts
for pupil area, so it is a direct measure of retinal light flux rather than scene brightness alone.

\section{Where it comes from}

Two steps, in \texttt{illuminance.py} and \texttt{pupil.py}:

\textbf{Step A --- the natural pupil diameter.} The Watson--Yellott (2012) unified pupil model,
solved iteratively because pupil size and retinal illuminance are mutually dependent (illuminance
depends on the effective stop, which depends on pupil size, which depends on illuminance). Converges
in a handful of iterations to 0.0001\,mm tolerance.

\textbf{Step B --- the effective stop.} Whichever is smaller, the pinhole or the natural pupil:
\begin{equation}
d_{\text{stop}} = \min(d_{\text{pupil}},\,d_{\text{pinhole}})
\end{equation}
For every one of your fabricated designs the pinhole is smaller (e.g.\ 0.400\,mm vs.\ a natural pupil
of 3.725\,mm at $M=8$), so \textbf{the pinhole is the effective aperture stop in every case} --- this
is the whole physical point of the device.

\begin{equation}
\text{Td} = L_{\text{scene}}\,[\text{cd/m}^2] \times A_{\text{stop}}\,[\text{mm}^2]
\end{equation}

\section{How it enters the acuity formula}

\begin{equation}
\Delta_{\text{bright}} = \begin{cases} 0 & \text{Td} \geq 300 \\[4pt]
0.060 \times \log_{10}(300/\text{Td}) & \text{Td} < 300 \end{cases}
\end{equation}

The 300\,Td knee is from Van Nes \& Bouman (1967): above it, contrast detection is governed by the
Weber--Fechner law and is illuminance-independent; below it, the de~Vries--Rose regime applies and
sensitivity falls off. The 0.060 logMAR-per-decade slope is from Tidbury, Czanner \& Newsham (2016),
verified against their Table~3 ($p<0.001$); the code comments record that an independently-derived
alternative slope from Sheedy, Bailey \& Raasch (1984), $\approx 0.0664$, could only be traced to a
secondary summary and so was not used.

\textbf{Worked example, 400\,\textmu m at $M=8$\,D:} $\text{Td} = 1256.6$, above the 300\,Td knee, so
$\Delta_{\text{bright}} = 0$.

\begin{warnbox}{The finding: this term is inert for every fabricated design, and you should say so first}
\begin{longtable}{ccc}
\textbf{Aperture (\textmu m)} & \textbf{Retinal illuminance (Td)} & \textbf{Brightness penalty
(logMAR)} \\
\midrule
75 & 44.2 & $+0.0499$ \\
100 & 78.5 & $+0.0349$ \\
150 & 176.7 & $+0.0138$ \\
\textbf{200} & \textbf{314.2} & \textbf{$+0.0000$} \\
\multicolumn{3}{c}{\emph{every fabricated design (400--2000\,\textmu m) falls here}} \\
2000 & 31\,415.9 & $+0.0000$ \\
\end{longtable}

The knee sits at 200\,\textmu m. Every one of your eight gaze-weighted winners is at 400\,\textmu m or
larger. At the model's assumed 10\,000\,cd/m$^2$ scene luminance, \textbf{the brightness term
contributes exactly zero logMAR to every design you have actually selected.} It is real physics,
correctly cited, correctly implemented --- and currently silent. It would activate for an indoor
measurement (typical luminance 100--300\,cd/m$^2$, roughly two decades below the assumption) or for
any future aperture below 200\,\textmu m. State this yourself: it pre-empts the question of whether
you actually tested the brightness term's effect on your headline result, and the honest answer is
that you have not needed to, because it does not act on that result.
\end{warnbox}

\chapter{What is computed but does NOT enter the recommended number}

Two further quantities are computed by the pipeline and reported alongside the recommended acuity
value, but neither contributes to it. Listing them matters because they appear in the same output
dictionary and could be mistaken for inputs if read from the code without care.

\section{MTF50}

The 50\%-contrast cutoff frequency, computed identically to MTF10 but at a higher contrast threshold.
Used only to bound the reported \texttt{approximate\_ci\_95\_mar\_arcmin} interval --- which, despite
the name, is \textbf{not a statistical confidence interval}. It is
$[\min, \max]$ of three different model choices (MTF10-based, MTF50-based, and the recommended
VSOTF-and-brightness-adjusted value), not a distribution with a defined confidence level.

\begin{warnbox}{Do not present this interval as a confidence interval}
At $M=8$\,D it spans logMAR $+0.318$ to $+0.795$ --- 4.8 chart lines. If shown, label it ``spread
across model choices,'' not ``95\% CI.'' An examiner asking what distribution produced a 95\% figure
has no defensible answer, because none was computed.
\end{warnbox}

\section{The legacy SNR term}

An older acuity model, retained in the code as \texttt{vsotf\_snr\_legacy} for direct comparison
rather than deleted. Uses an uncited $\text{SNR}^{0.05}$ exponent in place of the cited Tidbury
brightness correction. The project's own change log records that this was superseded specifically
because it lacked a literature source; it remains computable but is not the recommended value and
should not appear as a headline number in the presentation.

\chapter{Full input inventory}

Every value that can change the output of \texttt{predict\_acuity()}, traced to its origin.

\begin{longtable}{p{3.2cm}p{2.6cm}p{2.7cm}p{5.0cm}}
\caption{Complete input inventory for the acuity prediction.}\\
\toprule
\textbf{Input} & \textbf{Source} & \textbf{Design variable?} & \textbf{Reasoning / citation} \\
\midrule
\endhead
Aperture diameter $a_{\text{hole}}$ & user-set & \textbf{yes --- swept} & The physical design parameter
under optimisation \\
Wavelength & user-set, default 550\,nm & fixed & Photopic luminosity peak \\
Gap distance $d_{\text{gap}}$ & user-set, default 12\,mm & fixed & Spectacle vertex distance \\
Myopia $M$ & user-set & \textbf{yes --- swept} & The physiological variable under study \\
$P_0$, $n'$ & 60\,D, 1.336 & fixed & Standard emmetropic eye power and vitreous index \\
Scene luminance & 10\,000\,cd/m$^2$ & fixed default & Bright daylight reference; see
Ch.~3 warning on indoor measurement \\
Field diameter & 20$^\circ$ & fixed default & Typical CSF measurement field size \\
Subject age & 28.58\,yr & fixed default & Watson--Yellott reference age \\
Eyes (mono/bino) & 2 (binocular) & fixed default & Affects pupil-model flux density only \\
CSF model & Barten & fixed default & Most widely validated luminance-dependent CSF; two
alternatives computed but unused \\
Cycles/deg at 20/20 & 30.0 & fixed constant & Clinical definition, the entire scale anchor \\
VSOTF exponent & 0.1 & fixed constant & Deliberately weak; see Ch.~2 keybox \\
VSOTF clamp & $[0.25, 2.0]$ & fixed constant & Bounds the correction's maximum possible swing \\
Weber knee & 300\,Td & fixed constant, cited & Van Nes \& Bouman (1967), verified primary source \\
Brightness slope & 0.060 logMAR/decade & fixed constant, cited & Tidbury et al.\ (2016), verified
Table~3 \\
\bottomrule
\end{longtable}

\begin{keybox}{The reasoning behind the split between swept and fixed}
Only aperture and myopia vary across your actual results --- every other quantity in this table is a
fixed default representing a specific, stated measurement scenario (bright outdoor use, average adult
pupil dynamics, standard CSF assumptions). This is a reasonable scoping choice for a design
optimisation, but it means the predicted numbers are conditional on that scenario holding. Chapter~3's
finding --- that the brightness term is inert at this luminance --- is a direct consequence of that
choice, and is the one fixed default most likely to matter if the testing conditions differ from the
model's assumption.
\end{keybox}

\end{document}
